% Generated draft: requires mathematical and source review.
\documentclass[11pt]{article}
\usepackage[margin=25mm]{geometry}
\usepackage{fontspec}
\setmainfont{Latin Modern Roman}
\usepackage{amsmath,amssymb,mathtools,unicode-math}
\setmathfont{Latin Modern Math}
\usepackage{enumitem,needspace,xurl,hyperref}
\setlength{\parindent}{0pt}
\setlength{\parskip}{5pt}
\setlength{\emergencystretch}{3em}
\newcommand{\N}{\mathbb{N}}
\newcommand{\Z}{\mathbb{Z}}
\newcommand{\Q}{\mathbb{Q}}
\newcommand{\R}{\mathbb{R}}
\newcommand{\C}{\mathbb{C}}
\newcommand{\F}{\mathbb{F}}
\newcommand{\A}{\mathbb{A}}
\newcommand{\PP}{\mathbb{P}}
\newcommand{\E}{\mathbb{E}}
\newcommand{\eps}{\varepsilon}
\newcommand{\dd}{\,\mathrm{d}}
\begin{document}
\section*{UnsolvedMath ID: 800271}
\textbf{Source problem code:} catalogue entry
\subsection{Definitions and mathematical statement}
For each positive integer \(k\), let \(\PP^k(\C)\) denote complex projective \(k\)-space. If \(F:\PP^k(\C)\to\PP^k(\C)\) is holomorphic, write \(F^0=\operatorname{id}\) and \(F^{n+1}=F\circ F^n\) for integers \(n\geq 0\). The Fatou set \(\mathcal F(F)\) consists of points having an open neighborhood on which the family \(\{F^n:n\geq 0\}\) is normal: every sequence from this family has a subsequence converging uniformly on compact subsets, in the projective target, to a holomorphic map. A Fatou component is a nonempty connected component of \(\mathcal F(F)\).

For a given \(k\), the existence question is whether there are such a holomorphic map \(F\) and a Fatou component \(U\) satisfying
\[
F^n(U)\cap F^m(U)=\varnothing
\qquad\text{for all distinct integers }n,m\geq 0.
\]
Such a component is called wandering in the supplied statement.

The source leaves \(k\) unquantified, so it does not specify whether existence is requested in some positive dimension, in every positive dimension, or in a particular intended dimension. The question above is therefore presented as a family of questions indexed by \(k\), pending clarification. No restriction on the degree of \(F\) is supplied, and none is added here. Nonnegative iterate indices are the convention used in this draft; their range is not explicit in the supplied formula.
\subsection{Short English statement}
For complex projective space of dimension \(k\), ask whether a holomorphic self-map can have a Fatou component whose forward images are pairwise disjoint. The intended quantifier or restriction on \(k\) is unspecified, and no degree hypothesis is supplied. The registry labels the problem open; that label is supplied metadata, not an independently verified assessment.
\subsection{Sources}
\begin{itemize}
\item Local UnsolvedMath registry, \texttt{data/unsolvedmath/problems.json}, with additional records supplied through \texttt{data/unsolvedmath/additions/}: permanent ID \(800271\), exact source code \texttt{catalogue entry}, title \emph{Several Complex Variables Problem List --- Problem 70}. The supplied record lists Buzzard, Fornaess, and Krantz under attribution and gives the year 1995.
\item Buzzard, Fornaess, and Krantz, \emph{The Several Complex Variables Problem List} (1995), Problem 70. Supplied source URL: \url{https://arxiv.org/abs/math/9509202}. The supplied provenance records access on 2026-10-16 and extraction by \texttt{arxiv\_latex}; these are supplied provenance claims, not independently verified access or source checks.
\end{itemize}
\end{document}
